% Options for packages loaded elsewhere
\PassOptionsToPackage{unicode}{hyperref}
\PassOptionsToPackage{hyphens}{url}
\documentclass[
  japanese,
]{article}
\usepackage{xcolor}
\usepackage{amsmath,amssymb}
\setcounter{secnumdepth}{-\maxdimen} % remove section numbering
\usepackage{iftex}
\ifPDFTeX
  \usepackage[T1]{fontenc}
  \usepackage[utf8]{inputenc}
  \usepackage{textcomp} % provide euro and other symbols
\else % if luatex or xetex
  \usepackage{unicode-math} % this also loads fontspec
  \defaultfontfeatures{Scale=MatchLowercase}
  \defaultfontfeatures[\rmfamily]{Ligatures=TeX,Scale=1}
\fi
\usepackage{lmodern}
\ifPDFTeX\else
  % xetex/luatex font selection
\fi
% Use upquote if available, for straight quotes in verbatim environments
\IfFileExists{upquote.sty}{\usepackage{upquote}}{}
\IfFileExists{microtype.sty}{% use microtype if available
  \usepackage[]{microtype}
  \UseMicrotypeSet[protrusion]{basicmath} % disable protrusion for tt fonts
}{}
\makeatletter
\@ifundefined{KOMAClassName}{% if non-KOMA class
  \IfFileExists{parskip.sty}{%
    \usepackage{parskip}
  }{% else
    \setlength{\parindent}{0pt}
    \setlength{\parskip}{6pt plus 2pt minus 1pt}}
}{% if KOMA class
  \KOMAoptions{parskip=half}}
\makeatother
\ifLuaTeX
\usepackage[bidi=basic,shorthands=off,provide=*]{babel}
\else
\usepackage[bidi=default,shorthands=off,provide=*]{babel}
\fi
\ifLuaTeX
  \usepackage{selnolig} % disable illegal ligatures
\fi
\setlength{\emergencystretch}{3em} % prevent overfull lines
\providecommand{\tightlist}{%
  \setlength{\itemsep}{0pt}\setlength{\parskip}{0pt}}
\usepackage{bookmark}
\IfFileExists{xurl.sty}{\usepackage{xurl}}{} % add URL line breaks if available
\urlstyle{same}
% fallback for those not using the hyperref driver hyperxmp:
\makeatletter
\@ifundefined{xmpquote}{\newcommand{\xmpquote}[1]{#1}}{}
\makeatother
\hypersetup{
  pdftitle={イナールズィーラ語文法書},
  pdflang={ja},
  hidelinks,
  pdfcreator={LaTeX via pandoc}}

\title{イナールズィーラ語文法書}
\author{}
\date{}

\begin{document}
\maketitle

{
\setcounter{tocdepth}{3}
\tableofcontents
}
\section{\texorpdfstring{イナールズィーラ語文法書\\
{Cinarwziqili Gwlamwdufunw}
イナールズィーリ・グラームドゥフーン}{イナールズィーラ語文法書 Cinarwziqili Gwlamwdufunw イナールズィーリ・グラームドゥフーン}}\label{ux30a4ux30caux30fcux30ebux30baux30a3ux30fcux30e9ux8a9eux6587ux6cd5ux66f8-cinarwziqili-gwlamwdufunw-ux30a4ux30caux30fcux30ebux30baux30a3ux30fcux30eaux30b0ux30e9ux30fcux30e0ux30c9ux30a5ux30d5ux30fcux30f3}

\subsection{序章}\label{chapter1}

\subsubsection{はじめに}\label{section1.1}

\paragraph{この文法書について}\label{subsection1.1.1}

本書はイナールズィーラ帝国帝立大学言語学研究所によって刊行された辞典，{Cinarwziqili
Beruvwdikwtw} イナールズィーリ・ベルーヴディクト の巻末資料{Cinarwziqili
Gwlamwdufunw} イナールズィーリ・グラームドゥフーン
を日本語に翻訳したものである．

本辞典および文法書の翻訳にあたり，イナールズィーラ帝国帝立大学言語学研究所の皆様から多大なるご支援と貴重な助言を賜った．

ここに深く感謝の意を表する．

\subsubsection{言語の概要}\label{section1.2}

イナールズィーラ語は，砂漠の中のオアシス都市を中心に成立した帝国，イナールズィーラ帝国の公用語である．

その文法や発音は，五族共和連合時代に話されていたピスピ語から殆ど変化しておらず，多様な遺物が当時のまま解読可能である．

帝国議会によって制定された数多くの書体は，公開型電子書体として整備され，帝国が運営する電網頁にて公開されている．

\begin{center}\rule{0.5\linewidth}{0.5pt}\end{center}

\subsection{音韻体系}\label{chapter2}

\subsubsection{子音の発音}\label{section2.1}

子音文字は調音部位を象った形状を持ち，音声学的な特徴を視覚的に表現している．

\paragraph{子音発音解説}\label{subsection2.1.1}

\begin{description}
\tightlist
\item[{k}]
この文字は{/k/}という発音であることを指す．
\item[{g}]
この文字は{/g/}という発音であることを指す．
\item[{t}]
この文字は{/t/}という発音であることを指す．
\item[{d}]
この文字は{/d/}という発音であることを指す．
\item[{s}]
この文字は{/s/}という発音であることを指す．
\item[{z}]
この文字は{/z/}という発音であることを指す．
\item[{q}]
この文字は{/ʔ/}という発音であることを指す．
\item[{c}]
この文字は{/ʕ/}という発音であることを指す．
\item[{r}]
この文字は{/r/}という発音であることを指す．
\item[{l}]
この文字は{/l/}という発音であることを指す．
\item[{p}]
この文字は{/p/}という発音であることを指す．
\item[{b}]
この文字は{/b/}という発音であることを指す．
\item[{h}]
この文字は{/h/}という発音であることを指す．
\item[{x}]
この文字は{/x/}という発音であることを指す．
\item[{f}]
この文字は{/f/}という発音であることを指す．
\item[{v}]
この文字は{/v/}という発音であることを指す．
\item[{m}]
この文字は{/m/}という発音であることを指す．
\item[{n}]
この文字は{/n/}という発音であることを指す．
\end{description}

\paragraph{子音発音図表}\label{subsection2.1.2}

\subsubsection{母音の発音}\label{section2.2}

この言語において母音符号の付与は限定的であり，特に初学用や幼児用に母音符号を記述することもある．

\paragraph{母音発音解説}\label{subsection2.2.1}

\begin{description}
\tightlist
\item[{◌a}]
この文字は{/a/}という発音であることを指す．
\item[{◌e}]
この文字は{/e/}という発音であることを指す．
\item[{◌i}]
この文字は{/i/}という発音であることを指す．
\item[{◌o}]
この文字は{/o/}という発音であることを指す．
\item[{◌u}]
この文字は{/u/}という発音であることを指す．
\item[{◌w}]
この文字は無母音であることを指す．
\end{description}

\paragraph{母音発音図表}\label{subsection2.2.2}

\begin{center}\rule{0.5\linewidth}{0.5pt}\end{center}

\subsection{書記体系}\label{chapter3}

イナールズィーラ語の書記体系は，文字体系，標準字体，約物の三つの要素から構成される．

\subsubsection{文字体系}\label{section3.1}

\begin{description}
\tightlist
\item[ピスピ文字]
イナールズィーラ帝国成立以前の{Xwnasw Piswpi} ハナース・ピスピ
共和連合時代に広く用いられていた書記体系．現在では伝統的な装飾や詩歌などの筆記に用いられている．
\item[イナールズィーラ文字]
現行のイナールズィーラ語表記に広く用いられる文字．複数の字体があり，地域差が甚だしい書記体系とも取ることができるが，拡張文字も存在し，遍く使われ得る書記体系と銘打たれることもある．
\end{description}

\subsubsection{標準字体}\label{section3.2}

イナールズィーラ語の書記体系には，地域や時代によって異なる字体が存在する．これらの字体は，イナールズィーラ語の文化や歴史を反映している．

\paragraph{標準字体解説}\label{subsection3.2.1}

イナールズィーラ語に用いられる標準字体．

\begin{description}
\item[古代書体]
\begin{description}
\tightlist
\item[ピスピ体\\
{piswpi}]
五族共和連合時代に用いられていた書体．{Comwnwnaqanw} オムンナーン
という神が創造したと伝説上では語られている．
\item[スリーヴェ体\\
{sulive}]
五族共和連合時代に用いられていた書体である．{Sulive} スリーヴェ
というのは，ピスピ語で``流出させる''という意味を持つ動詞で，ピスピ体の筆記体として用いられている．
\end{description}
\item[中世書体]
\begin{description}
\tightlist
\item[コディート体\\
{kodito}]
イナールズィーラ帝国の工業革命期に使用された工業用規格書体である．後述のレクータ体の書記方向を判読性向上のために線形状にしている．
\item[レクータ体\\
{lekuta}]
\item[ロゼーグ体\\
{lozegw}]
\item[シルキ体\\
{silwki}]
\item[ハヴァーニ体\\
{xavani}]
イナールズィーラ帝国と東方の大帝国{Xavanw} ハヴァーン
の国交開通に伴って開発された書体である．何があっても揺るがないようにという意図が図形に込められている．
\item[ヘサーダ体\\
{xesada}]
イナールズィーラ帝国建国時にそれまで祀られていた五神{Comwnwnaqanw}
オムンナーン ，{Caxwkwsiqica} アフクスィーア ，{Pwratwluqusw}
プラートゥルース ，{Tanalapoteqe} タナーラポテー ，{Camiricaqalw}
アミーリアール に加えて，{Canumwnoqonw} アヌムノーン
という遍く中立の神を称え，六角形を象徴とした書体．
\item[ヒディーリ体\\
{xidili}]
イナールズィーラ帝国内での{Xavanw} ハヴァーン
文化の流行に即して制作された書体．{Xavanw} ハヴァーン
にて吉祥きっしょうを表す正八角形を基調としている．{xidili} ヒディーリ
とは{Xavanw} ハヴァーン にて吉祥きっしょうを表す言葉である．
\end{description}
\item[近世書体]
\begin{description}
\tightlist
\item[マキーナ体\\
{makina}]
イナールズィーラ帝国で最初期に用いられていたコンピュータ用フォントである．ビットマップフォントの一種で，低解像度のディスプレイに最適化されている．
\item[ポルゴ体\\
{polwgo}]
\item[ゾソーク体\\
{zosokw}]
イナールズィーラ帝国第五代目皇帝{Pwratuluqusw Cinalwgiqilw}
プラートゥルース・イナールズィール の皇后，{Camiricaqalw Zosokw}
アミーリアール・ゾソーク
が考案した書体．後にこれはイナールズィーラ帝国の教育用標準書体として制定された．
\end{description}
\end{description}

\subsubsection{約物}\label{section3.3}

イナールズィーラ語において，約物は主に文構造を明確にするために使用される．

\paragraph{約物解説}\label{subsection3.3.1}

イナールズィーラ語で用いられる約物の解説と役割を示す．

{.}

大休止符

{:}

強中休止符

{;}

弱中休止符

{,}

小休止符

{!}

感嘆符

{?}

疑問符

\begin{figure}
\centering
\caption{}\label{ltr-figure}
\end{figure}

右方符

\begin{figure}
\centering
\caption{}\label{rtl-figure}
\end{figure}

左方符

\begin{center}\rule{0.5\linewidth}{0.5pt}\end{center}

\subsection{子音の役割}\label{chapter4}

\subsubsection{単語構造}\label{section4.1}

\paragraph{基本単語構造}\label{subsection4.1.1}

イナールズィーラ語の基本単語構造はC\textsubscript{1}V\textsubscript{1}C\textsubscript{2}V\textsubscript{2}C\textsubscript{3}V\textsubscript{3}となる．これは原型となる単語のことであって，活用形などは後述する接辞の節にて単語構造が拡張される．

\paragraph{単語構造図表}\label{subsection4.1.2}

\subsubsection{子音概念}\label{section4.2}

語根はイナールズィーラ語において最も重要な単語の概念を指し示す水先案内人の様な存在である．後世の学者たちが単語を区別するために子音概念を明瞭化させた．そのために子音概念を理解することが重要である．

\paragraph{子音概念解説}\label{subsection4.2.1}

子音の概念に就いての説明

\begin{description}
\tightlist
\item[{k}]
この文字は剥離はくりの概念を指す．
\item[{g}]
この文字は癒着ゆちゃくの概念を指す．
\item[{t}]
この文字は乖離かいりの概念を指す．
\item[{d}]
この文字は同一どういつの概念を指す．
\item[{s}]
この文字は肉体にくたいの概念を指す．
\item[{z}]
この文字は精神せいしんの概念を指す．
\item[{q}]
この文字は空白くうはくの概念を指す．
\item[{c}]
この文字は物質ぶっしつの概念を指す．
\item[{r}]
この文字は過去かこの概念を指す．
\item[{l}]
この文字は未来みらいの概念を指す．
\item[{p}]
この文字は鎮静ちんせいの概念を指す．
\item[{b}]
この文字は高揚こうようの概念を指す．
\item[{h}]
この文字は受動じゅどうの概念を指す．
\item[{x}]
この文字は能動のうどうの概念を指す．
\item[{f}]
この文字は創造そうぞうの概念を指す．
\item[{v}]
この文字は破壊はかいの概念を指す．
\item[{m}]
この文字は流動りゅうどうの概念を指す．
\item[{n}]
この文字は固定こていの概念を指す．
\end{description}

\paragraph{子音概念図表}\label{subsection4.2.2}

\subsubsection{語根}\label{section4.3}

語根は単語の基本となる部分であり，語根に接辞を付けることで意味を変化させることができる．

\subsubsection{語根概念}\label{section4.4}

子音概念を三つ繋げることで語根概念が形成される．

\subsubsection{接頭辞}\label{section4.5}

接頭辞は語根の先頭に接続するものである．

\paragraph{接頭辞解説}\label{subsection4.5.1}

接頭辞に就いての説明

\begin{description}
\tightlist
\item[{k◌◌◌}]
この文字は接頭辞を剥離はくりの概念にさせる．
\item[{g◌◌◌}]
この文字は接頭辞を癒着ゆちゃくの概念にさせる．
\item[{t◌◌◌}]
この文字は接頭辞を乖離かいりの概念にさせる．
\item[{d◌◌◌}]
この文字は接頭辞を同一どういつの概念にさせる．
\item[{s◌◌◌}]
この文字は接頭辞を肉体にくたいの概念にさせる．
\item[{z◌◌◌}]
この文字は接頭辞を精神せいしんの概念にさせる．
\item[{q◌◌◌}]
この文字は接頭辞を空白くうはくの概念にさせる．
\item[{c◌◌◌}]
この文字は接頭辞を物質ぶっしつの概念にさせる．
\item[{r◌◌◌}]
この文字は接頭辞を過去かこの概念にさせる．
\item[{l◌◌◌}]
この文字は接頭辞を未来みらいの概念にさせる．
\item[{p◌◌◌}]
この文字は接頭辞を鎮静ちんせいの概念にさせる．
\item[{b◌◌◌}]
この文字は接頭辞を高揚こうようの概念にさせる．
\item[{h◌◌◌}]
この文字は接頭辞を受動じゅどうの概念にさせる．
\item[{x◌◌◌}]
この文字は接頭辞を能動のうどうの概念にさせる．
\item[{f◌◌◌}]
この文字は接頭辞を創造そうぞうの概念にさせる．
\item[{v◌◌◌}]
この文字は接頭辞を破壊はかいの概念にさせる．
\item[{m◌◌◌}]
この文字は接頭辞を流動りゅうどうの概念にさせる．
\item[{n◌◌◌}]
この文字は接頭辞を固定こていの概念にさせる．
\end{description}

\paragraph{接頭辞図表}\label{subsection4.5.2}

\subsubsection{接尾辞}\label{section4.6}

接尾辞は語根の末尾に接続するものである．

\paragraph{接尾辞解説}\label{subsection4.6.1}

接尾辞に就いての説明

\begin{description}
\tightlist
\item[{◌◌◌k}]
この文字は接尾辞を剥離はくりの概念にさせる．
\item[{◌◌◌g}]
この文字は接尾辞を癒着ゆちゃくの概念にさせる．
\item[{◌◌◌t}]
この文字は接尾辞を乖離かいりの概念にさせる．
\item[{◌◌◌d}]
この文字は接尾辞を同一どういつの概念にさせる．
\item[{◌◌◌s}]
この文字は接尾辞を肉体にくたいの概念にさせる．
\item[{◌◌◌z}]
この文字は接尾辞を精神せいしんの概念にさせる．
\item[{◌◌◌q}]
この文字は接尾辞を空白くうはくの概念にさせる．
\item[{◌◌◌c}]
この文字は接尾辞を物質ぶっしつの概念にさせる．
\item[{◌◌◌r}]
この文字は接尾辞を過去かこの概念にさせる．
\item[{◌◌◌l}]
この文字は接尾辞を未来みらいの概念にさせる．
\item[{◌◌◌p}]
この文字は接尾辞を鎮静ちんせいの概念にさせる．
\item[{◌◌◌b}]
この文字は接尾辞を高揚こうようの概念にさせる．
\item[{◌◌◌h}]
この文字は接尾辞を受動じゅどうの概念にさせる．
\item[{◌◌◌x}]
この文字は接尾辞を能動のうどうの概念にさせる．
\item[{◌◌◌f}]
この文字は接尾辞を創造そうぞうの概念にさせる．
\item[{◌◌◌v}]
この文字は接尾辞を破壊はかいの概念にさせる．
\item[{◌◌◌m}]
この文字は接尾辞を流動りゅうどうの概念にさせる．
\item[{◌◌◌n}]
この文字は接尾辞を固定こていの概念にさせる．
\end{description}

\paragraph{接尾辞図表}\label{subsection4.6.2}

\begin{center}\rule{0.5\linewidth}{0.5pt}\end{center}

\subsection{母音の役割}\label{chapter5}

イナールズィーラ語の格には否格，与格，属格，対格，主格，流格が存在する．

\subsubsection{前置格}\label{section5.1}

第一子音に附属する第一母音によって前置格が決定される．

\paragraph{前置格解説}\label{subsection5.1.1}

前置格に就いての説明

\begin{description}
\tightlist
\item[{◌a◌◌}]
否格，第一子音の概念を否定させる格．
\item[{◌e◌◌}]
与格，第一子音の概念を指向させる格．
\item[{◌i◌◌}]
属格，第一子音の概念を包含させる格．
\item[{◌o◌◌}]
対格，第一子音の概念を標的とする格．
\item[{◌u◌◌}]
主格，第一子音の概念を起点とする格．
\item[{◌w◌◌}]
流格，第一子音の概念を流出させる格．
\end{description}

\paragraph{前置格図表}\label{subsection5.1.2}

\subsubsection{後置格}\label{section5.2}

第二子音に附属する第二母音によって後置格が決定される．

\paragraph{後置格解説}\label{subsection5.2.1}

後置格に就いての説明

\begin{description}
\tightlist
\item[{◌◌a◌}]
否格，第二子音の概念を否定させる格．
\item[{◌◌e◌}]
与格，第二子音の概念を指向させる格．
\item[{◌◌i◌}]
属格，第二子音の概念を包含させる格．
\item[{◌◌o◌}]
対格，第二子音の概念を標的とする格．
\item[{◌◌u◌}]
主格，第二子音の概念を起点とする格．
\item[{◌◌w◌}]
流格，第二子音の概念を流出させる格．
\end{description}

\paragraph{後置格図表}\label{subsection5.2.2}

\subsubsection{複合格}\label{section5.3}

前置格と後置格の組み合わせによって複合格が決定される．

\paragraph{複合格解説}\label{subsection5.3.1}

複合格に就いての説明

\begin{description}
\item[{◌a◌◌}]
否格，第一子音の概念を否定させる格．

\begin{description}
\tightlist
\item[{◌a◌a◌}]
否否格，第一子音の概念を否定させ，第二子音の概念を否定させる格．
\item[{◌a◌e◌}]
否与格，第一子音の概念を否定させ，第二子音の概念を指向させる格．
\item[{◌a◌i◌}]
否属格，第一子音の概念を否定させ，第二子音の概念を包含させる格．
\item[{◌a◌o◌}]
否対格，第一子音の概念を否定させ，第二子音の概念を標的とする格．
\item[{◌a◌u◌}]
否主格，第一子音の概念を否定させ，第二子音の概念を起点とする格．
\item[{◌a◌w◌}]
否流格，第一子音の概念を否定させ，第二子音の概念を流出させる格．
\end{description}
\item[{◌e◌◌}]
与格，第一子音の概念を指向させる格．

\begin{description}
\tightlist
\item[{◌e◌a◌}]
与否格，第一子音の概念を指向させ，第二子音の概念を否定させる格．
\item[{◌e◌e◌}]
与与格，第一子音の概念を指向させ，第二子音の概念を指向させる格．
\item[{◌e◌i◌}]
与属格，第一子音の概念を指向させ，第二子音の概念を包含させる格．
\item[{◌e◌o◌}]
与対格，第一子音の概念を指向させ，第二子音の概念を標的とする格．
\item[{◌e◌u◌}]
与主格，第一子音の概念を指向させ，第二子音の概念を起点とする格．
\item[{◌e◌w◌}]
与流格，第一子音の概念を指向させ，第二子音の概念を流出させる格．
\end{description}
\item[{◌i◌◌}]
属格，第一子音の概念を所有する格．

\begin{description}
\tightlist
\item[{◌i◌a◌}]
属否格，第一子音の概念を包含させ，第二子音の概念を否定させる格．
\item[{◌i◌e◌}]
属与格，第一子音の概念を包含させ，第二子音の概念を指向させる格．
\item[{◌i◌i◌}]
属属格，第一子音の概念を包含させ，第二子音の概念を包含させる格．
\item[{◌i◌o◌}]
属対格，第一子音の概念を包含させ，第二子音の概念を標的とする格．
\item[{◌i◌u◌}]
属主格，第一子音の概念を包含させ，第二子音の概念を起点とする格．
\item[{◌i◌w◌}]
属流格，第一子音の概念を包含させ，第二子音の概念を流出させる格．
\end{description}
\item[{◌o◌◌}]
対格，第一子音の概念を対象とする格．

\begin{description}
\tightlist
\item[{◌o◌a◌}]
対否格，第一子音の概念を標的とし，第二子音の概念を否定させる格．
\item[{◌o◌e◌}]
対与格，第一子音の概念を標的とし，第二子音の概念を指向させる格．
\item[{◌o◌i◌}]
対属格，第一子音の概念を標的とし，第二子音の概念を包含させる格．
\item[{◌o◌o◌}]
対対格，第一子音の概念を標的とし，第二子音の概念を標的とする格．
\item[{◌o◌u◌}]
対主格，第一子音の概念を標的とし，第二子音の概念を起点とする格．
\item[{◌o◌w◌}]
対流格，第一子音の概念を標的とし，第二子音の概念を流出させる格．
\end{description}
\item[{◌u◌◌}]
主格，第一子音の概念を主語とする格．

\begin{description}
\tightlist
\item[{◌u◌a◌}]
主否格，第一子音の概念を主語とし，第二子音の概念を否定させる格．
\item[{◌u◌e◌}]
主与格，第一子音の概念を主語とし，第二子音の概念を指向させる格．
\item[{◌u◌i◌}]
主属格，第一子音の概念を主語とし，第二子音の概念を包含させる格．
\item[{◌u◌o◌}]
主対格，第一子音の概念を主語とし，第二子音の概念を標的とする格．
\item[{◌u◌u◌}]
主主格，第一子音の概念を主語とし，第二子音の概念を主語とする格．
\item[{◌u◌w◌}]
主流格，第一子音の概念を主語とし，第二子音の概念を流出させる格．
\end{description}
\item[{◌w◌◌}]
流格，第一子音の概念を流出させる格．

\begin{description}
\tightlist
\item[{◌w◌a◌}]
流否格，第一子音の概念を流出させ，第二子音の概念を否定させる格．
\item[{◌w◌e◌}]
流与格，第一子音の概念を流出させ，第二子音の概念を指向させる格．
\item[{◌w◌i◌}]
流属格，第一子音の概念を流出させ，第二子音の概念を包含させる格．
\item[{◌w◌o◌}]
流対格，第一子音の概念を流出させ，第二子音の概念を標的とする格．
\item[{◌w◌u◌}]
流主格，第一子音の概念を流出させ，第二子音の概念を起点とする格．
\item[{◌w◌w◌}]
流流格，第一子音の概念を流出させ，第二子音の概念を流出させる格．
\end{description}
\end{description}

\paragraph{複合格図表}\label{subsection5.3.2}

\subsubsection{接頭辞格}\label{section5.4}

接頭辞子音に附属する接頭辞母音によって接辞格が決定される．

\paragraph{接頭辞格解説}\label{subsection5.4.1}

接頭辞格に就いての説明

\begin{description}
\tightlist
\item[{◌a◌◌◌}]
否格，接頭辞子音の概念を否定させる格．
\item[{◌e◌◌◌}]
与格，接頭辞子音の概念を指向させる格．
\item[{◌i◌◌◌}]
属格，接頭辞子音の概念を包含させる格．
\item[{◌o◌◌◌}]
対格，接頭辞子音の概念を標的とする格．
\item[{◌u◌◌◌}]
主格，接頭辞子音の概念を起点とする格．
\item[{◌w◌◌◌}]
流格，接頭辞子音の概念を流出させる格．
\end{description}

\paragraph{接頭辞図表}\label{subsection5.4.2}

\subsubsection{接尾辞格}\label{section5.5}

接尾辞子音に附属する接尾辞母音によって接辞格が決定される．

\paragraph{接尾辞格解説}\label{subsection5.5.1}

接尾辞格に就いての説明

\begin{description}
\tightlist
\item[{◌◌◌◌a}]
否格，接尾辞子音の概念を否定させる格．
\item[{◌◌◌◌e}]
与格，接尾辞子音の概念を指向させる格．
\item[{◌◌◌◌i}]
属格，接尾辞子音の概念を包含させる格．
\item[{◌◌◌◌o}]
対格，接尾辞子音の概念を標的とする格．
\item[{◌◌◌◌u}]
主格，接尾辞子音の概念を起点とする格．
\item[{◌◌◌◌w}]
流格，接尾辞子音の概念を流出させる格．
\end{description}

\paragraph{接尾辞格図表}\label{subsection5.5.2}

\subsubsection{品詞}\label{section5.6}

イナールズィーラ語の品詞には附詞，動詞，容詞，助詞，副詞，名詞が存在する．

\paragraph{品詞解説}\label{subsection5.6.1}

品詞に就いての説明

\begin{description}
\tightlist
\item[{◌◌◌a}]
附詞，前置修飾して意味を限定させる品詞である．
\item[{◌◌◌e}]
動詞，動作の意味主体を表す品詞である．
\item[{◌◌◌i}]
容詞，物事の状態を表す品詞である．
\item[{◌◌◌o}]
助詞，後置修飾して意味を限定させる品詞である．
\item[{◌◌◌u}]
副詞，動作の意味周辺を表す品詞である．
\item[{◌◌◌w}]
名詞，物事の概念を表す品詞である．
\end{description}

\paragraph{品詞図表}\label{subsection5.6.2}

\begin{center}\rule{0.5\linewidth}{0.5pt}\end{center}

\subsection{文構造}\label{chapter6}

\subsubsection{基本文型}\label{section6.1}

主語，述語，目的語の配置

\begin{description}
\tightlist
\item[]
\end{description}

\subsubsection{複文と従属文}\label{section6.2}

複文の構成，従属節の使い方

\begin{description}
\tightlist
\item[]
\end{description}

\subsubsection{否定文と疑問文}\label{section6.3}

否定や疑問の表現方法

\begin{description}
\tightlist
\item[]
\end{description}

\subsubsection{強調と焦点}\label{section6.4}

文の中で特定の部分を強調する方法

\begin{description}
\tightlist
\item[]
\end{description}

\begin{center}\rule{0.5\linewidth}{0.5pt}\end{center}

\subsection{特殊構文}\label{chapter7}

\subsubsection{受動態}\label{section7.1}

受動態の形成と使用法

\begin{description}
\tightlist
\item[]
\end{description}

\subsubsection{能動態と中動態}\label{section7.2}

能動態，中動態の説明と使用例

\begin{description}
\tightlist
\item[]
\end{description}

\subsubsection{命令文}\label{section7.3}

命令文の形成と使用法

\begin{description}
\tightlist
\item[]
\end{description}

\begin{center}\rule{0.5\linewidth}{0.5pt}\end{center}

\subsection{辞書}\label{chapter8}

\subsubsection{辞書の使用方法}\label{section8.1}

イナールズィーラ語の辞書では主に語根からの派生形を捜査する形式となっている．

\subsubsection{新語彙の生成規則}\label{section8.2}

新しい語彙を作る規則

\begin{center}\rule{0.5\linewidth}{0.5pt}\end{center}

\subsection{付録}\label{chapter9}

\subsubsection{重要な語句}\label{section9.1}

イナールズィーラ語で頻出する語句を提示する．

\paragraph{頻出語句}\label{subsection9.1.1}

\begin{description}
\item[挨拶]
\begin{description}
\tightlist
\item[正式な挨拶]
{Kotiva!コティーヴァ！}
\item[くだけた挨拶]
{Xoqoci!ホーイ！}
\end{description}
\end{description}

\paragraph{神聖語句}\label{subsection9.1.2}

\begin{description}
\tightlist
\item[]
\end{description}

\begin{center}\rule{0.5\linewidth}{0.5pt}\end{center}

\subsection{活用形解説}\label{chapter10}

\subsubsection{附詞活用形}\label{section10.1}

\paragraph{附詞活用解説}\label{subsection10.1.1}

附詞の活用に就いての説明

\begin{description}
\item[{k◌◌◌a}]
剥離接頭辞

\begin{description}
\tightlist
\item[{k◌◌◌aq}]
剥離接頭辞空白接尾辞
\item[{k◌◌◌ac}]
剥離接頭辞物質接尾辞
\item[{k◌◌◌ar}]
剥離接頭辞過去接尾辞
\item[{k◌◌◌al}]
剥離接頭辞未来接尾辞
\item[{k◌◌◌ap}]
剥離接頭辞鎮静接尾辞
\item[{k◌◌◌ab}]
剥離接頭辞高揚接尾辞
\end{description}
\item[{g◌◌◌a}]
癒着接頭辞

\begin{description}
\tightlist
\item[{g◌◌◌aq}]
癒着接頭辞空白接尾辞
\item[{g◌◌◌ac}]
癒着接頭辞物質接尾辞
\item[{g◌◌◌ar}]
癒着接頭辞過去接尾辞
\item[{g◌◌◌al}]
癒着接頭辞未来接尾辞
\item[{g◌◌◌ap}]
癒着接頭辞鎮静接尾辞
\item[{g◌◌◌ab}]
癒着接頭辞高揚接尾辞
\end{description}
\item[{t◌◌◌a}]
乖離接頭辞

\begin{description}
\tightlist
\item[{t◌◌◌aq}]
乖離接頭辞空白接尾辞
\item[{t◌◌◌ac}]
乖離接頭辞物質接尾辞
\item[{t◌◌◌ar}]
乖離接頭辞過去接尾辞
\item[{t◌◌◌al}]
乖離接頭辞未来接尾辞
\item[{t◌◌◌ap}]
乖離接頭辞鎮静接尾辞
\item[{t◌◌◌ab}]
乖離接頭辞高揚接尾辞
\end{description}
\item[{d◌◌◌a}]
同一接頭辞

\begin{description}
\tightlist
\item[{d◌◌◌aq}]
同一接頭辞空白接尾辞
\item[{d◌◌◌ac}]
同一接頭辞物質接尾辞
\item[{d◌◌◌ar}]
同一接頭辞過去接尾辞
\item[{d◌◌◌al}]
同一接頭辞未来接尾辞
\item[{d◌◌◌ap}]
同一接頭辞鎮静接尾辞
\item[{d◌◌◌ab}]
同一接頭辞高揚接尾辞
\end{description}
\item[{s◌◌◌a}]
肉体接頭辞

\begin{description}
\tightlist
\item[{s◌◌◌aq}]
肉体接頭辞空白接尾辞
\item[{s◌◌◌ac}]
肉体接頭辞物質接尾辞
\item[{s◌◌◌ar}]
肉体接頭辞過去接尾辞
\item[{s◌◌◌al}]
肉体接頭辞未来接尾辞
\item[{s◌◌◌ap}]
肉体接頭辞鎮静接尾辞
\item[{s◌◌◌ab}]
肉体接頭辞高揚接尾辞
\end{description}
\item[{z◌◌◌a}]
精神接頭辞

\begin{description}
\tightlist
\item[{z◌◌◌aq}]
精神接頭辞空白接尾辞
\item[{z◌◌◌ac}]
精神接頭辞物質接尾辞
\item[{z◌◌◌ar}]
精神接頭辞過去接尾辞
\item[{z◌◌◌al}]
精神接頭辞未来接尾辞
\item[{z◌◌◌ap}]
精神接頭辞鎮静接尾辞
\item[{z◌◌◌ab}]
精神接頭辞高揚接尾辞
\end{description}
\end{description}

\paragraph{附詞活用図表}\label{subsection10.1.2}

附詞の格変化，数，定及び不定の概念など

\subsubsection{動詞活用形}\label{section10.2}

\paragraph{動詞活用解説}\label{subsection10.2.1}

動詞の活用に就いての説明

\begin{description}
\item[{h◌◌◌e}]
受動接頭辞

\begin{description}
\tightlist
\item[{h◌◌◌eq}]
受動接頭辞空白接尾辞
\item[{h◌◌◌ec}]
受動接頭辞物質接尾辞
\item[{h◌◌◌er}]
受動接頭辞過去接尾辞
\item[{h◌◌◌el}]
受動接頭辞未来接尾辞
\item[{h◌◌◌ep}]
受動接頭辞鎮静接尾辞
\item[{h◌◌◌eb}]
受動接頭辞高揚接尾辞
\end{description}
\item[{x◌◌◌e}]
能動接頭辞

\begin{description}
\tightlist
\item[{x◌◌◌eq}]
能動接頭辞空白接尾辞
\item[{x◌◌◌ec}]
能動接頭辞物質接尾辞
\item[{x◌◌◌er}]
能動接頭辞過去接尾辞
\item[{x◌◌◌el}]
能動接頭辞未来接尾辞
\item[{x◌◌◌ep}]
能動接頭辞鎮静接尾辞
\item[{x◌◌◌eb}]
能動接頭辞高揚接尾辞
\end{description}
\item[{f◌◌◌e}]
創造接頭辞

\begin{description}
\tightlist
\item[{f◌◌◌eq}]
創造接頭辞空白接尾辞
\item[{f◌◌◌ec}]
創造接頭辞物質接尾辞
\item[{f◌◌◌er}]
創造接頭辞過去接尾辞
\item[{f◌◌◌el}]
創造接頭辞未来接尾辞
\item[{f◌◌◌ep}]
創造接頭辞鎮静接尾辞
\item[{f◌◌◌eb}]
創造接頭辞高揚接尾辞
\end{description}
\item[{v◌◌◌e}]
破壊接頭辞

\begin{description}
\tightlist
\item[{v◌◌◌eq}]
破壊接頭辞空白接尾辞
\item[{v◌◌◌ec}]
破壊接頭辞物質接尾辞
\item[{v◌◌◌er}]
破壊接頭辞過去接尾辞
\item[{v◌◌◌el}]
破壊接頭辞未来接尾辞
\item[{v◌◌◌ep}]
破壊接頭辞鎮静接尾辞
\item[{v◌◌◌eb}]
破壊接頭辞高揚接尾辞
\end{description}
\item[{m◌◌◌e}]
流動接頭辞

\begin{description}
\tightlist
\item[{m◌◌◌eq}]
流動接頭辞空白接尾辞
\item[{m◌◌◌ec}]
流動接頭辞物質接尾辞
\item[{m◌◌◌er}]
流動接頭辞過去接尾辞
\item[{m◌◌◌el}]
流動接頭辞未来接尾辞
\item[{m◌◌◌ep}]
流動接頭辞鎮静接尾辞
\item[{m◌◌◌eb}]
流動接頭辞高揚接尾辞
\end{description}
\item[{n◌◌◌e}]
固定接頭辞

\begin{description}
\tightlist
\item[{n◌◌◌eq}]
固定接頭辞空白接尾辞
\item[{n◌◌◌ec}]
固定接頭辞物質接尾辞
\item[{n◌◌◌er}]
固定接頭辞過去接尾辞
\item[{n◌◌◌el}]
固定接頭辞未来接尾辞
\item[{n◌◌◌ep}]
固定接頭辞鎮静接尾辞
\item[{n◌◌◌eb}]
固定接頭辞高揚接尾辞
\end{description}
\end{description}

\paragraph{動詞活用図表}\label{subsection10.2.2}

動詞の活用，時制，態，相など

\subsubsection{容詞活用形}\label{section10.3}

\paragraph{容詞活用解説}\label{subsection10.3.1}

容詞の活用に就いての説明

\begin{description}
\item[{h◌◌◌i}]
受動接頭辞

\begin{description}
\tightlist
\item[{h◌◌◌ik}]
受動接頭辞剥離接尾辞
\item[{h◌◌◌ig}]
受動接頭辞癒着接尾辞
\item[{h◌◌◌it}]
受動接頭辞乖離接尾辞
\item[{h◌◌◌id}]
受動接頭辞同一接尾辞
\item[{h◌◌◌is}]
受動接頭辞肉体接尾辞
\item[{h◌◌◌iz}]
受動接頭辞精神接尾辞
\end{description}
\item[{x◌◌◌i}]
能動接頭辞

\begin{description}
\tightlist
\item[{x◌◌◌ik}]
能動接頭辞剥離接尾辞
\item[{x◌◌◌ig}]
能動接頭辞癒着接尾辞
\item[{x◌◌◌it}]
能動接頭辞乖離接尾辞
\item[{x◌◌◌id}]
能動接頭辞同一接尾辞
\item[{x◌◌◌is}]
能動接頭辞肉体接尾辞
\item[{x◌◌◌iz}]
能動接頭辞精神接尾辞
\end{description}
\item[{f◌◌◌i}]
創造接頭辞

\begin{description}
\tightlist
\item[{f◌◌◌ik}]
創造接頭辞剥離接尾辞
\item[{f◌◌◌ig}]
創造接頭辞癒着接尾辞
\item[{f◌◌◌it}]
創造接頭辞乖離接尾辞
\item[{f◌◌◌id}]
創造接頭辞同一接尾辞
\item[{f◌◌◌is}]
創造接頭辞肉体接尾辞
\item[{f◌◌◌iz}]
創造接頭辞精神接尾辞
\end{description}
\item[{v◌◌◌i}]
破壊接頭辞

\begin{description}
\tightlist
\item[{v◌◌◌ik}]
破壊接頭辞剥離接尾辞
\item[{v◌◌◌ig}]
破壊接頭辞癒着接尾辞
\item[{v◌◌◌it}]
破壊接頭辞乖離接尾辞
\item[{v◌◌◌id}]
破壊接頭辞同一接尾辞
\item[{v◌◌◌is}]
破壊接頭辞肉体接尾辞
\item[{v◌◌◌iz}]
破壊接頭辞精神接尾辞
\end{description}
\item[{m◌◌◌i}]
流動接頭辞

\begin{description}
\tightlist
\item[{m◌◌◌ik}]
流動接頭辞剥離接尾辞
\item[{m◌◌◌ig}]
流動接頭辞癒着接尾辞
\item[{m◌◌◌it}]
流動接頭辞乖離接尾辞
\item[{m◌◌◌id}]
流動接頭辞同一接尾辞
\item[{m◌◌◌is}]
流動接頭辞肉体接尾辞
\item[{m◌◌◌iz}]
流動接頭辞精神接尾辞
\end{description}
\item[{n◌◌◌i}]
固定接頭辞

\begin{description}
\tightlist
\item[{n◌◌◌ik}]
固定接頭辞剥離接尾辞
\item[{n◌◌◌ig}]
固定接頭辞癒着接尾辞
\item[{n◌◌◌it}]
固定接頭辞乖離接尾辞
\item[{n◌◌◌id}]
固定接頭辞同一接尾辞
\item[{n◌◌◌is}]
固定接頭辞肉体接尾辞
\item[{n◌◌◌iz}]
固定接頭辞精神接尾辞
\end{description}
\end{description}

\paragraph{容詞活用図表}\label{subsection10.3.2}

容詞の格変化，数，定及び不定の概念など

\subsubsection{助詞活用形}\label{section10.4}

\paragraph{助詞活用解説}\label{subsection10.4.1}

助詞の活用に就いての説明

\begin{description}
\item[{k◌◌◌o}]
剥離接頭辞

\begin{description}
\tightlist
\item[{k◌◌◌oq}]
剥離接頭辞空白接尾辞
\item[{k◌◌◌oc}]
剥離接頭辞物質接尾辞
\item[{k◌◌◌or}]
剥離接頭辞過去接尾辞
\item[{k◌◌◌ol}]
剥離接頭辞未来接尾辞
\item[{k◌◌◌op}]
剥離接頭辞鎮静接尾辞
\item[{k◌◌◌ob}]
剥離接頭辞高揚接尾辞
\end{description}
\item[{g◌◌◌o}]
癒着接頭辞

\begin{description}
\tightlist
\item[{g◌◌◌oq}]
癒着接頭辞空白接尾辞
\item[{g◌◌◌oc}]
癒着接頭辞物質接尾辞
\item[{g◌◌◌or}]
癒着接頭辞過去接尾辞
\item[{g◌◌◌ol}]
癒着接頭辞未来接尾辞
\item[{g◌◌◌op}]
癒着接頭辞鎮静接尾辞
\item[{g◌◌◌ob}]
癒着接頭辞高揚接尾辞
\end{description}
\item[{t◌◌◌o}]
乖離接頭辞

\begin{description}
\tightlist
\item[{t◌◌◌oq}]
乖離接頭辞空白接尾辞
\item[{t◌◌◌oc}]
乖離接頭辞物質接尾辞
\item[{t◌◌◌or}]
乖離接頭辞過去接尾辞
\item[{t◌◌◌ol}]
乖離接頭辞未来接尾辞
\item[{t◌◌◌op}]
乖離接頭辞鎮静接尾辞
\item[{t◌◌◌ob}]
乖離接頭辞高揚接尾辞
\end{description}
\item[{d◌◌◌o}]
同一接頭辞

\begin{description}
\tightlist
\item[{d◌◌◌oq}]
同一接頭辞空白接尾辞
\item[{d◌◌◌oc}]
同一接頭辞物質接尾辞
\item[{d◌◌◌or}]
同一接頭辞過去接尾辞
\item[{d◌◌◌ol}]
同一接頭辞未来接尾辞
\item[{d◌◌◌op}]
同一接頭辞鎮静接尾辞
\item[{d◌◌◌ob}]
同一接頭辞高揚接尾辞
\end{description}
\item[{s◌◌◌o}]
肉体接頭辞

\begin{description}
\tightlist
\item[{s◌◌◌oq}]
肉体接頭辞空白接尾辞
\item[{s◌◌◌oc}]
肉体接頭辞物質接尾辞
\item[{s◌◌◌or}]
肉体接頭辞過去接尾辞
\item[{s◌◌◌ol}]
肉体接頭辞未来接尾辞
\item[{s◌◌◌op}]
肉体接頭辞鎮静接尾辞
\item[{s◌◌◌ob}]
肉体接頭辞高揚接尾辞
\end{description}
\item[{z◌◌◌o}]
精神接頭辞

\begin{description}
\tightlist
\item[{z◌◌◌oq}]
精神接頭辞空白接尾辞
\item[{z◌◌◌oc}]
精神接頭辞物質接尾辞
\item[{z◌◌◌or}]
精神接頭辞過去接尾辞
\item[{z◌◌◌ol}]
精神接頭辞未来接尾辞
\item[{z◌◌◌op}]
精神接頭辞鎮静接尾辞
\item[{z◌◌◌ob}]
精神接頭辞高揚接尾辞
\end{description}
\end{description}

\paragraph{助詞活用図表}\label{subsection10.4.2}

助詞の機能と用法

\subsubsection{副詞活用形}\label{section10.5}

\paragraph{副詞活用解説}\label{subsection10.5.1}

副詞の活用に就いての説明

\begin{description}
\item[{h◌◌◌u}]
受動接頭辞

\begin{description}
\tightlist
\item[{h◌◌◌uq}]
受動接頭辞空白接尾辞
\item[{h◌◌◌uc}]
受動接頭辞物質接尾辞
\item[{h◌◌◌ur}]
受動接頭辞過去接尾辞
\item[{h◌◌◌ul}]
受動接頭辞未来接尾辞
\item[{h◌◌◌up}]
受動接頭辞鎮静接尾辞
\item[{h◌◌◌ub}]
受動接頭辞高揚接尾辞
\end{description}
\item[{x◌◌◌u}]
能動接頭辞

\begin{description}
\tightlist
\item[{x◌◌◌uq}]
能動接頭辞空白接尾辞
\item[{x◌◌◌uc}]
能動接頭辞物質接尾辞
\item[{x◌◌◌ur}]
能動接頭辞過去接尾辞
\item[{x◌◌◌ul}]
能動接頭辞未来接尾辞
\item[{x◌◌◌up}]
能動接頭辞鎮静接尾辞
\item[{x◌◌◌ub}]
能動接頭辞高揚接尾辞
\end{description}
\item[{f◌◌◌u}]
創造接頭辞

\begin{description}
\tightlist
\item[{f◌◌◌uq}]
創造接頭辞空白接尾辞
\item[{f◌◌◌uc}]
創造接頭辞物質接尾辞
\item[{f◌◌◌ur}]
創造接頭辞過去接尾辞
\item[{f◌◌◌ul}]
創造接頭辞未来接尾辞
\item[{f◌◌◌up}]
創造接頭辞鎮静接尾辞
\item[{f◌◌◌ub}]
創造接頭辞高揚接尾辞
\end{description}
\item[{v◌◌◌u}]
破壊接頭辞

\begin{description}
\tightlist
\item[{v◌◌◌uq}]
破壊接頭辞空白接尾辞
\item[{v◌◌◌uc}]
破壊接頭辞物質接尾辞
\item[{v◌◌◌ur}]
破壊接頭辞過去接尾辞
\item[{v◌◌◌ul}]
破壊接頭辞未来接尾辞
\item[{v◌◌◌up}]
破壊接頭辞鎮静接尾辞
\item[{v◌◌◌ub}]
破壊接頭辞高揚接尾辞
\end{description}
\item[{m◌◌◌u}]
流動接頭辞

\begin{description}
\tightlist
\item[{m◌◌◌uq}]
流動接頭辞空白接尾辞
\item[{m◌◌◌uc}]
流動接頭辞物質接尾辞
\item[{m◌◌◌ur}]
流動接頭辞過去接尾辞
\item[{m◌◌◌ul}]
流動接頭辞未来接尾辞
\item[{m◌◌◌up}]
流動接頭辞鎮静接尾辞
\item[{m◌◌◌ub}]
流動接頭辞高揚接尾辞
\end{description}
\item[{n◌◌◌u}]
固定接頭辞

\begin{description}
\tightlist
\item[{n◌◌◌uq}]
固定接頭辞空白接尾辞
\item[{n◌◌◌uc}]
固定接頭辞物質接尾辞
\item[{n◌◌◌ur}]
固定接頭辞過去接尾辞
\item[{n◌◌◌ul}]
固定接頭辞未来接尾辞
\item[{n◌◌◌up}]
固定接頭辞鎮静接尾辞
\item[{n◌◌◌ub}]
固定接頭辞高揚接尾辞
\end{description}
\end{description}

\paragraph{副詞活用図表}\label{subsection10.5.2}

副詞の格変化，数，定及び不定の概念など

\subsubsection{名詞活用形}\label{section10.6}

\paragraph{名詞活用解説}\label{subsection10.6.1}

名詞の活用に就いての説明

\begin{description}
\item[{h◌◌◌w}]
受動接頭辞

\begin{description}
\tightlist
\item[{h◌◌◌wk}]
受動接頭辞剥離接尾辞
\item[{h◌◌◌wg}]
受動接頭辞癒着接尾辞
\item[{h◌◌◌wt}]
受動接頭辞乖離接尾辞
\item[{h◌◌◌wd}]
受動接頭辞同一接尾辞
\item[{h◌◌◌ws}]
受動接頭辞肉体接尾辞
\item[{h◌◌◌wz}]
受動接頭辞精神接尾辞
\end{description}
\item[{x◌◌◌w}]
能動接頭辞

\begin{description}
\tightlist
\item[{x◌◌◌wk}]
能動接頭辞剥離接尾辞
\item[{x◌◌◌wg}]
能動接頭辞癒着接尾辞
\item[{x◌◌◌wt}]
能動接頭辞乖離接尾辞
\item[{x◌◌◌wd}]
能動接頭辞同一接尾辞
\item[{x◌◌◌ws}]
能動接頭辞肉体接尾辞
\item[{x◌◌◌wz}]
能動接頭辞精神接尾辞
\end{description}
\item[{f◌◌◌w}]
創造接頭辞

\begin{description}
\tightlist
\item[{f◌◌◌wk}]
創造接頭辞剥離接尾辞
\item[{f◌◌◌wg}]
創造接頭辞癒着接尾辞
\item[{f◌◌◌wt}]
創造接頭辞乖離接尾辞
\item[{f◌◌◌wd}]
創造接頭辞同一接尾辞
\item[{f◌◌◌ws}]
創造接頭辞肉体接尾辞
\item[{f◌◌◌wz}]
創造接頭辞精神接尾辞
\end{description}
\item[{v◌◌◌w}]
破壊接頭辞

\begin{description}
\tightlist
\item[{v◌◌◌wk}]
破壊接頭辞剥離接尾辞
\item[{v◌◌◌wg}]
破壊接頭辞癒着接尾辞
\item[{v◌◌◌wt}]
破壊接頭辞乖離接尾辞
\item[{v◌◌◌wd}]
破壊接頭辞同一接尾辞
\item[{v◌◌◌ws}]
破壊接頭辞肉体接尾辞
\item[{v◌◌◌wz}]
破壊接頭辞精神接尾辞
\end{description}
\item[{m◌◌◌w}]
流動接頭辞

\begin{description}
\tightlist
\item[{m◌◌◌wk}]
流動接頭辞剥離接尾辞
\item[{m◌◌◌wg}]
流動接頭辞癒着接尾辞
\item[{m◌◌◌wt}]
流動接頭辞乖離接尾辞
\item[{m◌◌◌wd}]
流動接頭辞同一接尾辞
\item[{m◌◌◌ws}]
流動接頭辞肉体接尾辞
\item[{m◌◌◌wz}]
流動接頭辞精神接尾辞
\end{description}
\item[{n◌◌◌w}]
固定接頭辞

\begin{description}
\tightlist
\item[{n◌◌◌wk}]
固定接頭辞剥離接尾辞
\item[{n◌◌◌wg}]
固定接頭辞癒着接尾辞
\item[{n◌◌◌wt}]
固定接頭辞乖離接尾辞
\item[{n◌◌◌wd}]
固定接頭辞同一接尾辞
\item[{n◌◌◌ws}]
固定接頭辞肉体接尾辞
\item[{n◌◌◌wz}]
固定接頭辞精神接尾辞
\end{description}
\end{description}

\paragraph{名詞活用図表}\label{subsection10.6.2}

名詞の格変化，数，定及び不定の概念など

\begin{center}\rule{0.5\linewidth}{0.5pt}\end{center}

\end{document}
